% Generated by scripts/assemble.py; corrections belong in source/corrections.json.
\examhead{50}{上海财经大学 2024 年数学分析试卷}{51}
% p052-b09; source PDF page 52
\par\Needspace{5\baselineskip}
\blockanchor{p052-b09}
\majorquestion{一}解答题.\par
% p052-b10; source PDF page 52
\blockanchor{p052-b10}
\question{1}计算 $\lim_{x\to +\infty}\left(x^{\frac{2}{x}} - 1\right)^{\frac{2}{\ln x}}$\par
% p052-b11; source PDF page 52
\blockanchor{p052-b11}
\question{2}计算 $\lim_{n\to \infty}\int_0^1\frac{(2n - 1)x^{n - 1}}{1 + 2^x}\mathrm{d}x.$\par
% p052-b12; source PDF page 52
\blockanchor{p052-b12}
\question{3}设 $f(x) = \begin{cases} \sin \frac{x}{2}, & x > 0; \\ x \arctan x, & x \leq 0. \end{cases}$ 求 $\int_{0}^{\pi + 1} f(x - 1) \, \mathrm{d}x$.\par
% p052-b13; source PDF page 52
\blockanchor{p052-b13}
\question{4}解答如下问题:\par
% p052-b14; source PDF page 52
\blockanchor{p052-b14}
(1) 计算 $f(x) = \frac{1}{1 - x - x^2}$ 在 $x = 0$ 处的幂级数展开式.\par
% p052-b15; source PDF page 52
\blockanchor{p052-b15}
(2) 设 $a_0 = a_1 = 1$, $a_n = a_{n-2} + a_{n-1} (n \geq 2)$, 求幂级数 $\sum_{n=0}^{\infty} a_n x^n$ 的和函数.\par
% p052-b16; source PDF page 52
\blockanchor{p052-b16}
\question{5}求一曲线旋转得到的旋转体的体积.\par
% p052-b17; source PDF page 52
\blockanchor{p052-b17}
\question{6}求曲线积分 $\int_{L} y^{3} \mathrm{~d} x + |3x^{2} - y| \mathrm{~d} y,$ 其中 $L: x^{2} + y^{2} = 4,$ 取逆时针方向.\par
% p052-b18; source PDF page 52
\par\Needspace{6\baselineskip}
\blockanchor{p052-b18}
\question{7}求曲面积分\par
\nopagebreak[4]
% p052-b19; source PDF page 52
\blockanchor{p052-b19}
\displayformula{\iint_ {\Sigma} \left(y ^ {2} - z ^ {2}\right) \mathrm{d} y \mathrm{d} z + \left(z ^ {2} - x ^ {2}\right) \mathrm{d} z \mathrm{d} x + \left(x ^ {2} - y ^ {2}\right) \mathrm{d} x \mathrm{d} y.}
% p052-b20; source PDF page 52
\blockanchor{p052-b20}
其中 $\Sigma$ 为上半椭球面 $\frac{x^2}{a^2} +\frac{y^2}{b^2} +\frac{z^2}{c^2} = 1,z\geq 0(a,b,c > 0)$ ，方向取上侧.\par
% p052-b21; source PDF page 52
\blockanchor{p052-b21}
\question{8}已知  $\left(\frac{\partial f}{\partial x}\right)^{2} + \left(\frac{\partial f}{\partial y}\right)^{2} = 4,$  作变换  $\left\{\begin{aligned}x &= uv \\ y &= \frac{1}{2}(u^{2} - v^{2})\end{aligned}\right.$  使得  $f(x,y)$  变为  $g(u,v)$ .\par
% p052-b22; source PDF page 52
\blockanchor{p052-b22}
(1) 若 $a\left(\frac{\partial g}{\partial u}\right)^2 + b\left(\frac{\partial g}{\partial v}\right)^2 = u^2 + v^2$，求 $a, b$.\par
% p052-b23; source PDF page 52
\blockanchor{p052-b23}
(2) 对 (1) 式求出的 $a, b$, 利用极坐标变换 $u = r \cos \theta, v = r \sin \theta$ 将 $g$ 转化为 $r, \theta$ 的表达式.\par
% p052-b24; source PDF page 52
\par\Needspace{5\baselineskip}
\blockanchor{p052-b24}
\majorquestion{二}证明题.\par
% p053-b00; source PDF page 53
\blockanchor{p053-b00}
\question{1}讨论 $\int_0^1\frac{\sin^2x}{x^p(x^p + \sin x)}\mathrm{d}x$ 在 $p$ 取什么范围时收敛.\par
% p053-b01; source PDF page 53
\blockanchor{p053-b01}
\question{2}解答如下问题:\par
% p053-b02; source PDF page 53
\blockanchor{p053-b02}
(1) 设 $u_{n}(x)$ 是 $[a, b]$ 上的单调递增函数, 且 $\sum_{n=1}^{\infty} u_{n}(a)$ 与 $\sum_{n=1}^{\infty} u_{n}(b)$ 绝对收敛, 问 $\sum_{n=1}^{\infty} u_{n}(x)$ 在 $[a, b]$ 上是否绝对收敛? 是否一致收敛?\par
% p053-b03; source PDF page 53
\blockanchor{p053-b03}
(2) 设 $u_{n}(x)$ 是 $[a, b]$ 上的单调递增函数, 且 $\sum_{n=1}^{\infty} u_{n}(a)$ 与 $\sum_{n=1}^{\infty} u_{n}(b)$ 收敛, 问 $\sum_{n=1}^{\infty} u_{n}(x)$ 在 $[a, b]$ 上是一致收敛?\par
% p053-b04; source PDF page 53
\blockanchor{p053-b04}
\question{3}设函数 $f(x)$ 在 $[0,1]$ 上二阶可导，且 $f(0) = f(1) = 0$，$\min_{0 \leq x \leq 1} f(x) = -1$，证明：$\max_{0 \leq x \leq 1} f''(x) \geq 8$.\par
% p053-b05; source PDF page 53
\par\Needspace{6\baselineskip}
\blockanchor{p053-b05}
\question{4}设 $f(x)$ 在 $[0,1]$ 上恒正连续, $m, M$ 是 $f(x)$ 在此区间上的最小值和最大值, 证明:\par
\nopagebreak[4]
% p053-b06; source PDF page 53
\blockanchor{p053-b06}
\displayformula{1 \leq \int_ {0} ^ {1} \frac {1}{f (x)} \mathrm{d} x \int_ {0} ^ {1} f (x) \mathrm{d} x \leq \frac {(m + M) ^ {2}}{4 M m}.}
% p053-b07; source PDF page 53
\par\Needspace{6\baselineskip}
\blockanchor{p053-b07}
\question{5}设函数 $f(x)$ 满足对任意的 $x \in (a, b)$, 存在 $L_x > 0$ 及 $\delta_x > 0$, 满足\par
\nopagebreak[4]
% p053-b08; source PDF page 53
\blockanchor{p053-b08}
\displayformula{\left| f \left(x _ {1}\right) - f \left(x _ {2}\right) \right| <   L _ {x} \left| x _ {1} - x _ {2} \right|, \forall x _ {1}, x _ {2} \in \left(x - \delta_ {x}, x + \delta_ {x}\right) \cap (a, b).}
% p053-b09; source PDF page 53
\blockanchor{p053-b09}
证明: 存在 L > 0, 对任意的  $x, y \in (a, b)$ , 成立  $|f(x) - f(y)| < L|x - y|$ .\par
% p053-b10; source PDF page 53
\blockanchor{p053-b10}
\question{6}证明 $\int_{1}^{+\infty} \frac{\cos(xt)}{1 + x^2} \, \mathrm{d}x$ 在 $(- \infty, +\infty)$ 上一致收敛, 并证明 $\lim_{t \to +\infty} \int_{1}^{+\infty} \frac{\cos(xt)}{1 + x^2} \, \mathrm{d}x = 0$.\par
